% Guia técnico do UnDF-TeX.
% Compilação: LuaLaTeX, em ciclos até estabilizar o sumário.

\DocumentMetadata{lang=pt-BR,tagging=on}
\tagpdfsetup{math/mathml/luamml/load=false}
\documentclass[12pt,a4paper,oneside]{report}
\usepackage{../estilos/manual-style-guia}

\hypersetup{
  pdftitle={Guia de Uso do UnDF-TeX},
  pdfauthor={Universidade do Distrito Federal Professor Jorge Amaury Maia Nunes},
  pdfsubject={Documentação técnica do UnDF-TeX 0.1.0},
  pdfkeywords={UnDF, UnDF-TeX, LuaLaTeX, BibLaTeX, Biber, versão 1.0.0}
}

\begin{document}

\pagenumbering{gobble}

\begin{titlepage}
  \centering
  \vspace*{1.5cm}
  {\large Universidade do Distrito Federal Professor Jorge Amaury Maia Nunes -- UnDF\par}
  \vspace{3cm}
  {\bfseries\Large GUIA DE USO DO UnDF-TeX\par}
  \vspace{0.45cm}
  {\bfseries\normalsize PARTE II\par}
  \vspace{0.8cm}
  {\normalsize Preparação, edição, compilação e verificação técnica do template\par}
  \vfill
  {\normalsize Versão do guia: 1.0.0\par}
  {\normalsize Compatível com o UnDF-TeX: 0.1.0\par}
  {\normalsize Situação: minuta para validação institucional\par}
  {\normalsize Brasília\par}
  {\normalsize 20 de outubro de 2026\par}
\end{titlepage}

\manualfrontmatter

\chapter*{Apresentação}
\addcontentsline{toc}{chapter}{Apresentação}

Este guia documenta o uso do UnDF-TeX, desde a preparação do ambiente até a geração e a conferência do PDF. Seu público inclui estudantes que adotam o template, pessoas que prestam apoio técnico e responsáveis por manutenção.

O guia explica implementação e operação. As regras acadêmicas, metodológicas e de normalização pertencem ao \emph{Manual de Trabalhos de Conclusão de Curso da UnDF}. Quando um comando produz capa, citação, tabela ou outro elemento, este documento descreve como usar o recurso; o manual acadêmico explica o que o elemento significa e quando deve ser utilizado.

\manualwarning{O UnDF-TeX 0.1.0 é um projeto técnico em validação. Sua existência não torna o uso obrigatório e não substitui PPC, regulamento de curso, orientação, decisão de colegiado ou ato institucional.}

\clearpage
\tableofcontents

\manualmainmatter

\chapter{Apresentação da Ferramenta}

\section{O que é o UnDF-TeX}

O UnDF-TeX é uma classe e um conjunto de arquivos LaTeX destinados à preparação de Trabalhos de Conclusão de Curso. O projeto separa conteúdo, metadados, elementos opcionais, bibliografia e implementação interna para reduzir formatação manual e tornar a saída reproduzível.

O template gera capa, folha de rosto, folha de aprovação, resumos, sumário, capítulos, referências e elementos opcionais a partir dos dados e arquivos preenchidos pelo usuário. Também oferece comandos para citações, figuras, tabelas, fluxogramas, equações, pseudocódigos e referências cruzadas.

\section{Finalidade e limites}

O objetivo é automatizar tarefas repetitivas e detectar problemas técnicos comuns. O template não avalia mérito acadêmico, não verifica a veracidade das fontes, não decide se uma pesquisa precisa de autorização ética e não substitui revisão humana.

\manualnote{Automação de aparência não transforma uma escolha do projeto em regra da ABNT nem em decisão institucional. O manual acadêmico e as orientações do curso continuam sendo as referências para conteúdo e aplicabilidade.}

\section{Versões e situação}

Este guia corresponde à revisão documental \texttt{1.0.0} e descreve o UnDF-TeX \texttt{1.0.0}. Antes de iniciar ou atualizar um trabalho, confira o arquivo \texttt{VERSION.txt}, o \href{https://github.com/igormrqs/undftex-publico/blob/main/CHANGELOG.md}{histórico de alterações} e a origem do pacote recebido. A versão do manual e a versão técnica do template têm controles próprios.

Versões do projeto podem mudar. Preserve uma cópia funcional do trabalho antes de substituir a classe ou os módulos internos e faça nova regressão do PDF depois da atualização.

\section{Relação com o manual acadêmico}

Este guia remete ao Manual de TCC sempre que o usuário precisa decidir se um elemento é obrigatório, opcional ou adequado. Aqui são documentados apenas localização, sintaxe, processamento, mensagens e comportamento técnico.

\chapter{Preparação do Ambiente}

\section{Requisitos}

O projeto exige uma distribuição LaTeX atualizada com LuaLaTeX e Biber. Também é necessário um editor de texto com suporte a arquivos UTF-8 e, para conferência, um leitor de PDF.

Ambientes locais podem usar MiKTeX ou TeX Live. Em serviços online, selecione LuaLaTeX como compilador e envie a estrutura completa do projeto, preservando nomes e caminhos relativos.

\section{Obtenção da versão correta}

Obtenha o pacote pelo repositório controlado do projeto ou por canal institucional indicado. Não copie apenas \texttt{tcc.tex}: a classe, os módulos de \texttt{src/}, a identidade visual e os arquivos bibliográficos são necessários.

Depois de extrair ou clonar o projeto, confirme a presença de:

\begin{lstlisting}
undftex.cls
tcc.tex
src/
identidade/
obrigatorio/
opcional/
capitulos/
\end{lstlisting}

\section{Teste inicial}

Compile o modelo em branco antes de inserir conteúdo. O teste separa problemas de instalação de erros introduzidos durante a edição.

\begin{lstlisting}
lualatex -interaction=nonstopmode -halt-on-error tcc.tex
biber tcc
lualatex -interaction=nonstopmode -halt-on-error tcc.tex
lualatex -interaction=nonstopmode -halt-on-error tcc.tex
\end{lstlisting}

O resultado esperado é \texttt{tcc.pdf}. Mensagens sobre campos ainda não preenchidos são normais no primeiro teste; interrupção do processo, arquivo ausente ou comando desconhecido exigem correção.

\section{Editor e codificação}

Mantenha os arquivos em UTF-8. Não substitua caracteres acentuados por sequências incompatíveis nem altere a codificação de apenas um capítulo. Configure o editor para mostrar final de linha, pares de chaves e erros de sintaxe.

\chapter{Organização dos Arquivos}

\section{Entradas públicas}

A entrada da monografia é \texttt{tcc.tex}. Para consultar recursos demonstrativos,
use \texttt{tcc-exemplo-undf.tex} e a pasta \texttt{exemplo/}. Não desenvolva o
trabalho diretamente sobre o exemplo completo, pois ele contém elementos
opcionais e textos destinados a teste.

Na branch de desenvolvimento, as demais entradas são \texttt{tcc-artigo.tex}
(TCC-artigo), \texttt{artigo.tex} (artigo autônomo),
\texttt{relatorio-projeto-aplicado.tex} (relatório parcial de cada projeto),
\texttt{portfolio.tex} (relatório final do portfólio, modalidade 1.1),
\texttt{relatorio-projeto.tex}
(relatório final de projeto técnico/profissional, modalidade 1.3) e
\texttt{negocio-tecnologico.tex} (relatório final de negócio tecnológico,
modalidade 1.4). O índice \texttt{modelos/README.md} relaciona cada modelo ao
seu exemplo e aos arquivos editáveis. A base pública possui desenvolvimento
independente e não inclui automaticamente essas modalidades experimentais.

O relatório parcial preserva memória, conteúdo e evidências de cada projeto.
Sua abertura é enxuta, sem capa de TCC, aprovação, resumos ou sumário; não
exige periodicidade, parecer ou cronograma. Edite a identificação, o relato
e as fontes consultadas em \texttt{modelos/projeto-aplicado/}.

O portfólio final reúne todos os projetos do percurso declarado: introdução,
organização, inventário, síntese de cada projeto, análise integradora e
conclusão. Não basta concatenar os registros. Edite os metadados, resumos,
conteúdo e bibliografia em \texttt{modelos/portfolio/}. O exemplo preenchido
está em \texttt{portfolio-exemplo-undf.tex}, com arquivos em
\texttt{exemplo/portfolio/}. O número de projetos do exemplo é demonstrativo.

Os elementos obrigatórios e opcionais e as referências do portfólio final
preservam o padrão do TCC convencional. Artigo científico não é exigido
nesses modelos nem no relatório profissional. O documento de agosto de 2026
foi referência de organização do relato conclusivo. A admissibilidade e o
recorte do percurso devem ser confirmados com o curso.

\section{Mapa do projeto}

\begin{longtable}{p{0.31\linewidth}p{0.57\linewidth}}
\toprule
Caminho & Finalidade \\
\midrule
\endfirsthead
\toprule
Caminho & Finalidade \\
\midrule
\endhead
\texttt{tcc.tex} & Entrada do modelo em branco e ordem de montagem do documento. \\
\texttt{capitulos/} & Conteúdo textual da monografia. \\
\texttt{obrigatorio/} & Resumos e base bibliográfica da monografia. \\
\texttt{artigo/} & Metadados, resumos e miolo dos modelos de artigo. \\
\texttt{modelos/} & Índice e arquivos preenchíveis de relatório parcial, portfólio, projeto profissional e negócio tecnológico; apresentação comum dos TCCs finais. \\
\texttt{exemplo/} & Conteúdo didático independente dos modelos para preenchimento. \\
\texttt{opcional/} & Errata, dedicatória, agradecimentos, epígrafe, glossário, apêndices e anexos. \\
\texttt{figuras/} & Imagens usadas pelo trabalho. \\
\texttt{identidade/} & Ativos institucionais carregados pelo template. \\
\texttt{src/core/} & Metadados, estrutura, títulos, layout, arquivos e referências cruzadas. \\
\texttt{src/abnt/} & Camadas técnicas organizadas por norma. \\
\texttt{src/identity/} & Capa e páginas de identificação. \\
\texttt{src/biblatex/} & Estilo e comandos bibliográficos. \\
\bottomrule
\end{longtable}

\section{Arquivos que o usuário edita}

Edite \texttt{tcc.tex}, os arquivos de \texttt{capitulos/}, \texttt{obrigatorio/}, \texttt{opcional/} e os ativos inseridos em \texttt{figuras/}. Arquivos internos de \texttt{src/} devem ser modificados apenas durante manutenção do template e acompanhados de testes nos modelos em branco e de exemplo.

Para as demais modalidades, edite a entrada escolhida e a pasta indicada no
índice dos modelos. Mantenha os arquivos de exemplos separados do trabalho real.
No TCC, os dados catalográficos continuam obrigatórios na versão final; o caminho
\texttt{opcional/ficha-catalografica.tex} é mantido por compatibilidade e não
define a obrigatoriedade do elemento.

\section{Inclusão de capítulos}

O arquivo principal usa \texttt{\textbackslash undfincludechapter}. Para ativar a seção opcional de dados ou corpus, retire o comentário da linha correspondente:

\begin{lstlisting}
\undfincludechapter{capitulos/04-dados}
\end{lstlisting}

Para adicionar outra seção primária, crie o arquivo na pasta \texttt{capitulos/} e inclua-o na posição desejada. Evite renomear arquivos internos apenas por preferência estética; caminhos inconsistentes são causa frequente de falha.

\section{Arquivos auxiliares}

LuaLaTeX e Biber criam arquivos como \texttt{.aux}, \texttt{.bcf}, \texttt{.bbl}, \texttt{.log}, \texttt{.toc}, \texttt{.lof} e \texttt{.lot}. Eles não devem ser editados manualmente. Use limpeza controlada quando houver estado inconsistente e recompile o documento completo.

\chapter{Informações Gerais do Trabalho}

As informações gerais, também chamadas de metadados, ficam no preâmbulo de \texttt{tcc.tex}. Elas alimentam páginas iniciais e propriedades do PDF.

\section{Autoria, título e vínculo}

\begin{lstlisting}
\autor{Nome do Aluno}
% \autores{Aluno 1//Aluno 2}
\titulo{Título do Trabalho}
\subtitulo{Subtítulo opcional}
\curso{Nome do curso}
\tipotrabalho{Trabalho de Conclusão de Curso}
\graupretendido{bacharel em Nome do curso}
\local{Cidade}
\ano{Ano}
\end{lstlisting}

Use \texttt{\textbackslash autor} para uma autoria. Para múltiplas autorias admitidas pelo curso, use \texttt{\textbackslash autores} e separe nomes com \texttt{//}. Não use os dois comandos ao mesmo tempo.

\section{Orientação e aprovação}

\begin{lstlisting}
\orientador{Nome do Orientador}
\coorientador{}
\areadeconcentracao{Área de concentração}
\dataaprovacao{Brasília, DF, dia de mês de ano}
\end{lstlisting}

Deixe o argumento de \texttt{\textbackslash coorientador} vazio quando não houver coorientação. A data e o formato da banca devem acompanhar a decisão institucional aplicável; os valores do modelo são exemplos editáveis.

\section{Banca examinadora}

Limpe a lista padrão e cadastre cada membro com nome, prefixo acadêmico e vínculo:

\begin{lstlisting}
\limparmembrosbanca
\membrobanca{Nome do Orientador}{Prof. Dr.}{UnDF}
\membrobanca{Nome do Membro 1}{Profa. Dra.}{Instituição}
\membrobanca{Nome do Membro 2}{Prof. Dr.}{Instituição}
\end{lstlisting}

O primeiro argumento contém o nome sem titulação. O segundo é impresso antes do nome. O terceiro aparece na linha de vínculo institucional. Cadastre apenas a composição efetivamente aprovada.

\section{Conferência}

Depois de preencher os campos, confira capa, folha de rosto, folha de aprovação e propriedades do PDF. Um erro no preâmbulo pode aparecer em mais de uma página.

\chapter{Elementos Pré-textuais no \marcaUnDFTeX}

Este capítulo descreve a implementação. Consulte a Parte I para obrigatoriedade, conteúdo e finalidade acadêmica.

\section{Montagem automática}

O arquivo principal chama, em ordem, capa, folha de rosto, verso, errata, folha de aprovação, elementos opcionais, resumos e sumário. Não duplique esses comandos dentro dos arquivos de conteúdo.

\begin{lstlisting}
\undfcover
\pretextual
\undffolhaderosto
\undfversofolhaderostoopcional
\undferrata
\undffolhadeaprovacao
\end{lstlisting}

\section{Resumos}

Preencha \texttt{obrigatorio/resumo.tex}:

\begin{lstlisting}
\undfsetcontent{resumo}{Texto do resumo.}
\undfsetcontent{palavraschave}{termo 1; termo 2; termo 3}
\end{lstlisting}

No arquivo \texttt{obrigatorio/abstract.tex}, informe título, texto e palavras-chave. Para idioma diferente do inglês, altere também o rótulo:

\begin{lstlisting}
\undfsetcontent{abstracttitle}{RESUMEN}
\undfsetcontent{abstract}{Texto em língua estrangeira.}
\undfsetcontent{keywordslabel}{Palabras clave}
\undfsetcontent{keywords}{término 1; término 2}
\end{lstlisting}

O template acrescenta o ponto ao final das palavras-chave. Por isso, informe os termos separados por ponto e vírgula, sem ponto adicional no argumento.

\section{Dados catalográficos}

O arquivo \texttt{opcional/ficha-catalografica.tex} aceita conteúdo textual, imagem ou PDF conforme os comandos públicos do modelo. Utilize apenas o material fornecido pelo fluxo institucional competente.

\begin{lstlisting}
\dadoscatalograficos{Texto fornecido pela unidade responsável.}
% \fichacatalograficaimagem{figuras/ficha.png}
% \fichacatalograficapdf{figuras/ficha.pdf}
\end{lstlisting}

\section{Errata}

\begin{lstlisting}
\errataref{Referência completa do trabalho.}
\erratatexto{%
  \erratatabela{%
    \erratalinha{16}{10}{auto-clavado}{autoclavado}
  }
}
\end{lstlisting}

\section{Elementos opcionais}

Arquivos vazios ou inexistentes não são impressos. O carregamento é feito pelo comando mostrado a seguir. Preencha apenas os elementos usados e não remova a lógica condicional do arquivo principal sem necessidade.

\begin{lstlisting}
\undfinputoptional{opcional/nome-do-arquivo}
\end{lstlisting}

\section{Listas e sumário}

\begin{lstlisting}
\listadeilustracoes
\listadefluxogramas
\listadetabelas
\undfsumario
\end{lstlisting}

O modelo em branco chama apenas o sumário por padrão. Ative listas quando forem adequadas ao trabalho e recompile até que títulos e páginas se estabilizem.

\chapter{Citações, Notas e Referências}

Este capítulo descreve a operação dos recursos bibliográficos. Os critérios para selecionar fontes, citar com integridade e elaborar referências são tratados na Parte I.

\section{Arquivo bibliográfico}

Cadastre as obras em \texttt{obrigatorio/referencias.bib} e associe o arquivo ao documento principal:

\begin{lstlisting}
\addbibresource{obrigatorio/referencias.bib}
\end{lstlisting}

Cada entrada possui uma chave única, usada pelos comandos de citação. Evite espaços, acentos e chaves duplicadas. Sempre confira no PDF se nomes, título, edição, data, endereço eletrônico e data de acesso foram processados como esperado.

\section{Citações no texto}

\begin{lstlisting}
\citarparenteses{chave}
\citarparenteses[p. 25]{chave}
\citartexto{chave}
\citartexto[p. 25]{chave}

\citardireta[p. 25]{chave}{Texto citado.}
\citardiretatexto[p. 25]{chave}{Texto citado.}
\end{lstlisting}

O argumento entre colchetes é opcional e recebe o localizador, por exemplo \texttt{p. 25}. O argumento seguinte é a chave da entrada bibliográfica.

Para citação direta longa, use:

\begin{lstlisting}
\begin{citacaolonga}[p. 25]{chave}
Texto da citação.
\end{citacaolonga}
\end{lstlisting}

Para citação de citação, prefira informar autoria e ano da original como texto,
usando uma chave apenas para a fonte efetivamente consultada:

\begin{lstlisting}
\parenciteapud{Freire}{1994}[13]{consultada}[25]
\textapud{Freire}{1994}[13]{consultada}[25]
\end{lstlisting}

Os comandos anteriores \texttt{\string\apudbib} e \texttt{\string\apudbibonline}
continuam aceitando duas chaves. Ao usar \texttt{\string\printundfreferences},
a obra original usada somente nessa citação indireta é omitida da lista; uma
citação independente da original na mesma seção bibliográfica a inclui. Essa
filtragem não é aplicada a chamadas diretas a \texttt{\string\printbibliography}.

\section{Notas de rodapé}

\begin{lstlisting}
\notarodape{Texto explicativo da nota.}
\end{lstlisting}

\section{Lista de referências}

Insira \texttt{\string\printundfreferences} após a parte textual. A lista é produzida apenas com as entradas citadas, salvo configuração bibliográfica adicional.

\section{Diagnóstico bibliográfico}

Uma citação com a chave impressa em negrito ou como marcador indefinido geralmente indica chave inexistente ou ciclo de compilação incompleto. Mensagens como \texttt{Please (re)run Biber} pedem a execução do Biber e novas passagens do LuaLaTeX.

\chapter{Figuras, Tabelas, Equações e Outros Objetos}

\section{Figura individual}

\begin{lstlisting}
\undffigura[0.75\textwidth]
  {figuras/exemplo.jpg}
  {Legenda da figura}
  {rotulo}
  [Autor, ano]
\end{lstlisting}

A largura e a fonte são opcionais. Sem indicação, o comando usa \texttt{0.8\string\textwidth} e \texttt{elaboração própria}. Referencie com \texttt{\string\figuraref\{rotulo\}}.

\section{Duas figuras independentes lado a lado}

\begin{lstlisting}
\undffigurasladoalado
  {figuras/antes.jpg}{Registro anterior}{antes}[elaboração própria]
  {figuras/depois.jpg}{Registro posterior}{depois}[elaboração própria]
\end{lstlisting}

O comando cria duas figuras independentes, cada uma com número, legenda, rótulo e fonte. Não há legenda coletiva nem subdivisões ``a'' e ``b''. A largura opcional inicial controla cada coluna.

\section{Tabela}

\begin{lstlisting}
\begin{undftabela}{Legenda da tabela}{rotulo}[Fonte]
  \begin{tabular}{ll}
    Item & Valor \\
  \end{tabular}
\end{undftabela}
\end{lstlisting}

Use \texttt{\string\tabelaref\{rotulo\}} para a chamada no texto. Tabelas extensas que atravessam páginas devem ser construídas com um ambiente multipágina apropriado, como \texttt{longtable}, pois \texttt{undftabela} é um flutuante de uma página.

\section{Gráficos e quadros}

Na versão 0.1.0, um gráfico fornecido como imagem é inserido por \texttt{\string\undffigura}; a saída o identifica na categoria geral ``Figura''. Um gráfico produzido diretamente em LaTeX pode ser colocado em um ambiente \texttt{figure}, com \texttt{\string\caption}, rótulo iniciado por \texttt{fig:} e \texttt{\string\undffonte}.

Não há ambiente público dedicado a quadro na versão 0.1.0. Não use \texttt{undftabela} apenas para obter bordas quando o objeto for semanticamente um quadro, pois a identificação sairia como ``Tabela''. Até que uma interface própria seja incorporada e testada, quadros exigem implementação manual controlada pelo mantenedor. Essa limitação está registrada para uma versão posterior e não deve ser ocultada por documentação de um comando inexistente.

\section{Equações}

\begin{lstlisting}
\equacao[energia]{E = mc^2}
\equacaosemnumero{x + y = z}
\end{lstlisting}

Somente a equação numerada e rotulada pode ser referenciada:

\begin{lstlisting}
\equacaoref{energia}
\end{lstlisting}

\section{Fluxograma}

\begin{lstlisting}
\begin{undffluxograma}{Fluxo de análise}{analise}
  \centering
  \begin{tikzpicture}
    % nós e conexões
  \end{tikzpicture}
\end{undffluxograma}
\end{lstlisting}

Até três ocorrências, o recurso integra a numeração e a lista de figuras. Acima desse total, passa a usar numeração e lista próprias. Após adicionar ou remover fluxogramas, compile novamente até os rótulos e as listas estabilizarem.

\section{Pseudocódigo}

\begin{lstlisting}
\begin{undfpseudocodigo}{MEDIA(V)}{media}{Calcular a média}
Entrada: lista V
Saida: média dos valores
...
\end{undfpseudocodigo}
\end{lstlisting}

O ambiente comum tenta iniciar o bloco apenas quando há espaço suficiente para cerca de meia página. Para um algoritmo que precise continuar na página seguinte, use \texttt{undfpseudocodigolongo}; ele permite a quebra interna do bloco. Referencie ambos com \texttt{\string\pseudocodigoref\{rotulo\}}.

\section{Apêndices, anexos e PDFs externos}

\begin{lstlisting}
\undfapendice{Título}{rotulo}
\undfanexo{Título}{rotulo}
\undfapendicepdf{Título}{arquivo.pdf}
\undfanexarpdf{Título}{arquivo.pdf}
\end{lstlisting}

Para inclusão sem criar apêndice ou anexo, use \texttt{\string\undfincluirpdf\{arquivo.pdf\}}. Os comandos de referência são \texttt{\string\apendiceref} e \texttt{\string\anexoref}.

\section{Imagens e paginação}

Os comandos de figura verificam o tamanho do arquivo e emitem avisos para imagens pesadas. Uma figura fixada que não caiba no espaço restante avança inteira para a próxima página, podendo deixar área em branco. Isso preserva a unidade do objeto. Ajuste a largura, reposicione a chamada no texto ou use uma imagem mais compacta; não force sobreposição com texto ou margens.

\chapter{Compilação e Geração do PDF}

\section{Fluxo automatizado}

Na raiz do projeto, execute:

\begin{lstlisting}
latexmk -outdir=build/perfis
\end{lstlisting}

O arquivo \texttt{.latexmkrc} seleciona LuaLaTeX, aciona o Biber quando necessário e define \texttt{tcc.tex} como entrada padrão. Para o exemplo completo:

\begin{lstlisting}
latexmk -outdir=build/perfis tcc-exemplo-undf.tex
\end{lstlisting}

\section{Fluxo manual}

Quando o \texttt{latexmk} não estiver disponível, use:

\begin{lstlisting}
New-Item -ItemType Directory -Force build/perfis | Out-Null
lualatex -output-directory=build/perfis tcc.tex
biber --input-directory=build/perfis --output-directory=build/perfis tcc
lualatex -output-directory=build/perfis tcc.tex
lualatex -output-directory=build/perfis tcc.tex
\end{lstlisting}

No PowerShell, as saídas ficam em \texttt{build/perfis/}. O script
\texttt{scripts/verificar-perfis.ps1} automatiza esse ciclo no Windows, inclusive
sem Perl. O Biber deve receber o nome-base, sem a extensão \texttt{.tex}.
Passagens adicionais podem ser necessárias após alterações em sumário, listas,
citações ou referências cruzadas.

\section{Overleaf}

Envie a estrutura completa do projeto e selecione LuaLaTeX como compilador. O arquivo \texttt{.latexmkrc} acompanha o repositório. Confirme no painel de logs se a bibliografia foi processada e se não restaram referências indefinidas.

\section{Arquivos auxiliares}

Arquivos como \texttt{.aux}, \texttt{.bcf}, \texttt{.bbl}, \texttt{.toc}, \texttt{.lof}, \texttt{.lot} e \texttt{.lofl} registram dados entre as passagens. Eles não fazem parte do texto autoral e podem ser regenerados. Para limpar artefatos pelo \texttt{latexmk}, use \texttt{latexmk -outdir=build/perfis -c}; faça nova compilação completa em seguida.

\section{Critério de conclusão}

Considere a geração tecnicamente estabilizada somente quando: o comando terminar sem erro fatal; o PDF abrir normalmente; não houver pedido de nova execução do Biber; sumário, listas e referências cruzadas estiverem resolvidos; e a inspeção visual das páginas críticas tiver sido feita.

\chapter{Validações e Avisos}

\section{Estimativa de ocupação dos elementos pré-textuais}

O arquivo principal chama, antes da capa:

\begin{lstlisting}
\undfvalidateonepagecontents
\end{lstlisting}

O comando estima o espaço ocupado por resumo, resumo em língua estrangeira, dedicatória, agradecimentos, epígrafe, errata e ficha catalográfica textual. A medição é preventiva: caixas, fontes, imagens e outros elementos podem fazer o resultado visual diferir da estimativa. Sempre confira o PDF.

A estimativa é um recurso técnico do projeto e não cria, por si só, obrigação normativa ou institucional de limitar o elemento a uma página. Quando houver limite definido pelo curso ou por decisão institucional aplicável, essa origem deve ser identificada separadamente.

\section{Resumo e resumo em língua estrangeira}

Por padrão, excesso estimado de conteúdo gera aviso técnico no log e a compilação continua. Para transformar esse caso em erro intencional durante uma revisão controlada, use antes da validação:

\begin{lstlisting}
\bloquearresumoslongos
\undfvalidateonepagecontents
\end{lstlisting}

Para restabelecer explicitamente o modo de aviso:

\begin{lstlisting}
\avisarresumoslongos
\end{lstlisting}

O sistema também estima a contagem de palavras e avisa quando o resumo parece ficar fora da faixa configurada. Essa contagem pode divergir de editores de texto diante de comandos, siglas, números e palavras compostas.

O modo de erro é uma ferramenta optativa de revisão e não representa o comportamento oficial de entrega. Uma recomendação normativa ou editorial, isoladamente, não justifica bloquear a compilação.

\section{Conteúdo ausente}

\texttt{\string\undfwarnifempty} produz aviso quando um conteúdo obrigatório não foi preenchido. Arquivos opcionais vazios ou inexistentes são ignorados por \texttt{\string\undfinputoptional}.

\section{Imagens}

Os comandos de figura podem avisar sobre arquivos acima dos limiares configurados. Para suspender e reativar a verificação:

\begin{lstlisting}
\undfdesabilitaravisosimagem
\undfhabilitaravisosimagem
\end{lstlisting}

Esses avisos avaliam peso do arquivo, não nitidez, contraste, legibilidade nem adequação de resolução. A conferência visual continua obrigatória.

\section{Natureza dos diagnósticos}

\begin{description}
  \item[Erro] interrompe ou invalida a geração; deve ser corrigido.
  \item[Aviso] sinaliza risco ou pendência; exige avaliação humana.
  \item[Informação] registra uma decisão do processamento e não implica defeito.
  \item[Estimativa] aproxima ocupação ou contagem e sempre exige conferência humana.
\end{description}

Requisito normativo, decisão institucional validada, recomendação e regra de curso pertencem à documentação acadêmica; aviso, estimativa e erro de compilação pertencem ao diagnóstico técnico. Não trate um PDF gerado como prova automática de conformidade. O template automatiza forma e diagnósticos limitados, mas não avalia conteúdo, método, autoria ou decisão institucional.

\chapter{Solução de Problemas}

\begin{longtable}{>{\raggedright\arraybackslash}p{0.25\textwidth} >{\raggedright\arraybackslash}p{0.31\textwidth} >{\raggedright\arraybackslash}p{0.34\textwidth}}
\toprule
\textbf{Sintoma} & \textbf{Causa provável} & \textbf{Ação recomendada} \\
\midrule
\endhead
Classe não encontrada & Compilação fora da raiz ou pacote incompleto & Compile na raiz e confira \texttt{undftex.cls} e a pasta \texttt{src/}. \\
\addlinespace
LuaLaTeX exigido & Motor incorreto & Selecione LuaLaTeX no editor ou serviço. \\
\addlinespace
Citações indefinidas & Chave ausente ou Biber não executado & Confira o \texttt{.bib}, execute Biber e repita LuaLaTeX. \\
\addlinespace
Referência cruzada com ``??'' & Passagens insuficientes ou rótulo divergente & Confira o rótulo e recompile até estabilizar. \\
\addlinespace
Sumário ou listas desatualizados & Arquivos auxiliares antigos & Faça novas passagens; se persistir, limpe auxiliares e recompile. \\
\addlinespace
Figura não encontrada & Caminho, extensão ou maiúsculas divergentes & Confira o nome real e use caminho relativo à raiz. \\
\addlinespace
Fonte ausente ou substituída & Distribuição incompleta ou cache de fontes & Atualize a distribuição, confira o log e não aceite substituição silenciosa sem revisar o PDF. \\
\addlinespace
Caractere especial causa erro & Símbolo reservado ou arquivo fora de UTF-8 & Preserve UTF-8 e use a sintaxe LaTeX adequada para caracteres reservados. \\
\addlinespace
Grande espaço antes de figura & Objeto fixado não cabe no restante da página & Reduza a largura, reposicione a chamada ou aceite o avanço integral. \\
\addlinespace
Página em branco inesperada & Quebra explícita, inclusão vazia ou ajuste de objeto & Localize a transição no fonte e verifique \texttt{clearpage}, arquivos opcionais e objetos anteriores. \\
\addlinespace
Conteúdo fora da margem & Linha de código, URL, tabela ou imagem larga & Reduza ou permita quebra; procure \texttt{Overfull} no log e confira a página. \\
\addlinespace
Pseudocódigo ocupa página seguinte & Bloco não cabe ou é longo & Use \texttt{undfpseudocodigolongo} quando a quebra interna for necessária. \\
\addlinespace
Fluxograma com lista incorreta & Contagem mudou desde a última passagem & Recompile até desaparecer o aviso de contagem. \\
\addlinespace
Resumo gera aviso de uma página & Estimativa de altura excedida & Revise o texto e confira o PDF; o padrão não bloqueia a compilação. \\
\addlinespace
PDF externo com página cortada & Dimensões ou opções de inclusão incompatíveis & Ajuste opções de \texttt{pdfpages} no argumento opcional e confira todas as páginas. \\
\addlinespace
Compilação interrompida & Erro de sintaxe, arquivo ausente ou ferramenta indisponível & Corrija a primeira mensagem iniciada por \texttt{!} e execute novamente o ciclo completo. \\
\bottomrule
\end{longtable}

\section{Leitura do log}

Procure primeiro a primeira linha iniciada por \texttt{!}, pois erros posteriores podem ser consequência dela. Em seguida, busque \texttt{Undefined}, \texttt{Warning}, \texttt{Overfull} e pedidos de nova compilação. Registre o comando executado e a primeira mensagem relevante ao solicitar suporte.

\section{Isolamento do problema}

Quando a causa não estiver clara, reproduza o erro no menor conjunto possível: comente temporariamente a inclusão do capítulo suspeito, reduza o conteúdo até identificar o trecho e restaure as partes gradualmente. Faça esse diagnóstico em cópia ou sob controle de versão para não perder conteúdo.

\chapter{Checklist Técnico de Fechamento}

\section{Arquivos e versão}

\begin{itemize}[label=$\square$]
  \item Confirme que o trabalho usa uma cópia íntegra do UnDF-TeX 0.1.0.
  \item Confira os arquivos de entrada, capítulos, imagens e bibliografia.
  \item Remova apenas artefatos temporários; preserve fontes e arquivos necessários à reprodução.
  \item Verifique se não há dados de exemplo remanescentes.
\end{itemize}

\section{Compilação}

\begin{itemize}[label=$\square$]
  \item Compile com LuaLaTeX e Biber pelo fluxo recomendado.
  \item Repita as passagens até sumário, listas e rótulos estabilizarem.
  \item Corrija todos os erros fatais e avalie cada aviso relevante.
  \item Confirme que não há citações ou referências cruzadas indefinidas.
\end{itemize}

\section{Conferência visual}

\begin{itemize}[label=$\square$]
  \item Abra o PDF gerado, não apenas a visualização do editor.
  \item Confira capa, folha de rosto, aprovação, resumos e sumário.
  \item Examine transições de página, objetos extensos e páginas com pouco conteúdo.
  \item Confira nitidez, legenda, fonte e chamada de cada objeto.
  \item Verifique links, marcadores, numeração e páginas incluídas externamente.
  \item Confirme ordem de leitura, contraste, legibilidade e alternativas textuais quando o fluxo de produção as suportar.
  \item Faça uma leitura em escala de 100\% e outra em visão geral de miniaturas.
\end{itemize}

\section{Entrega técnica}

\begin{itemize}[label=$\square$]
  \item Gere o PDF a partir da versão efetivamente aprovada para entrega.
  \item Use nome de arquivo claro, sem marcas como \texttt{final-final}.
  \item Confirme abertura em outro leitor de PDF.
  \item Arquive o código-fonte, a bibliografia e os recursos usados na geração.
\end{itemize}

\chapter{Compatibilidade, Licença e Manutenção}

\section{Matriz de compatibilidade}

\begin{longtable}{p{0.27\textwidth}p{0.25\textwidth}p{0.38\textwidth}}
\toprule
\textbf{Documento ou componente} & \textbf{Versão} & \textbf{Situação} \\
\midrule
\endhead
Guia de Uso do UnDF-TeX & 1.0.0 & Compatível com a versão técnica indicada nesta edição. \\
Manual de TCC da UnDF & 1.0.0 & Documento acadêmico autônomo relacionado. \\
UnDF-TeX & 0.1.0 & Versão de referência deste guia. \\
LuaLaTeX e Biber & Distribuições atuais compatíveis & Devem oferecer os pacotes carregados pelo projeto. \\
\bottomrule
\end{longtable}

\section{Versionamento}

Mudanças na classe ou nos pacotes podem alterar comandos e aparência. Antes de atualizar um trabalho em andamento, registre a versão atual, crie uma cópia ou uma ramificação e compare os PDFs antes e depois. Não misture arquivos internos de versões diferentes.

Consulte o \href{https://github.com/igormrqs/undftex-publico/blob/main/CHANGELOG.md}{histórico público de alterações}
e a \href{https://github.com/igormrqs/undftex-publico}{página do repositório}
para identificar a origem do pacote. As versões editoriais 1.0.0 do manual
e do guia não representam, por si só, atualização do template para 1.0.0.

\section{Licença}

A LPPL 1.3c ou posterior aplica-se aos arquivos de código expressamente identificados. Consulte \texttt{LICENSE} e \texttt{LICENSE-LPPL-1.3c} antes de modificar ou redistribuir esses arquivos.

Marcas, logotipos, elementos de identidade institucional, fontes, imagens, documentos oficiais e recursos de terceiros permanecem sujeitos às respectivas regras de uso e não são automaticamente licenciados pela LPPL. A documentação textual somente deve ser redistribuída sob a licença que vier a ser expressamente indicada para ela.

\section{Responsabilidades de manutenção}

Uma alteração técnica deve atualizar código, exemplo, testes pertinentes, documentação de uso, histórico e indicação de compatibilidade. A validação deve incluir compilação do modelo em branco e do exemplo completo, comparação visual e revisão dos avisos do log.

\section{Relato de problema}

Informe versão, ambiente, entrada, comando, primeira mensagem e trecho mínimo. Em defeitos visuais, inclua página ou captura; omita dados pessoais e conteúdo sigiloso.

Problemas relativos à base pública podem ser registrados na
\href{https://github.com/igormrqs/undftex-publico/issues}{área de issues do repositório}.
Para recursos experimentais, use o canal de acompanhamento disponibilizado
pelo responsável pelo desenvolvimento privado.


\appendix
\chapter{Referência Rápida de Comandos}

As tabelas reúnem a interface pública mais usada. Nos exemplos, argumentos entre colchetes são opcionais; argumentos entre chaves são obrigatórios.

\section{Informações gerais}

\begin{longtable}{>{\ttfamily\footnotesize\raggedright\arraybackslash}p{0.40\textwidth}p{0.20\textwidth}p{0.30\textwidth}}
\toprule
\normalfont\textbf{Comando} & \textbf{Finalidade} & \textbf{Uso e observações} \\
\midrule
\endhead
\string\autor & Autor único & Parâmetro: nome. Use \string\autores para trabalho coletivo. \\
\string\autores & Dois ou mais autores & Parâmetro: nomes separados por \texttt{//}. \\
\string\titulo & Título principal & Parâmetro: texto sem a pontuação do subtítulo. \\
\string\subtitulo & Subtítulo opcional & Parâmetro: texto; deixe vazio quando inexistente. \\
\string\curso & Curso & Parâmetro: nome por extenso. \\
\string\tipotrabalho & Tipo documental & Parâmetro: tipo do trabalho. \\
\string\graupretendido & Grau pretendido & Parâmetro: texto que completa a natureza. \\
\string\local & Local & Parâmetro: cidade. \\
\string\ano & Ano & Parâmetro: quatro algarismos. \\
\string\orientador & Orientador & Parâmetro: nome sem prefixo acadêmico. \\
\string\coorientador & Coorientador & Parâmetro: nome; pode ficar vazio. \\
\string\areadeconcentracao & Área de concentração & Parâmetro: área, quando pertinente. \\
\string\dataaprovacao & Data de aprovação & Parâmetro: local e data por extenso. \\
\string\membrobanca & Membro da banca & Parâmetros: nome, prefixo e instituição. \\
\bottomrule
\end{longtable}

\section{Conteúdo e arquivos}

\begin{longtable}{>{\ttfamily\footnotesize\raggedright\arraybackslash}p{0.40\textwidth}p{0.20\textwidth}p{0.30\textwidth}}
\toprule
\normalfont\textbf{Comando} & \textbf{Finalidade} & \textbf{Uso e observações} \\
\midrule
\endhead
\string\undfsetcontent & Registrar conteúdo & Parâmetros: chave e conteúdo. \\
\string\undfusecontent & Recuperar conteúdo & Parâmetro: chave. \\
\string\undfifcontent & Testar conteúdo & Parâmetros: chave, saída preenchida e saída vazia. \\
\string\undfinputoptional & Incluir arquivo opcional & Parâmetro: caminho sem \texttt{.tex}. \\
\string\undfincludechapter & Incluir capítulo & Parâmetro: caminho sem \texttt{.tex}. \\
\string\undfwarnifempty & Avisar ausência & Parâmetros: chave e mensagem; não interrompe. \\
\string\undfincluirpdf & Incluir PDF integral & Parâmetro: arquivo; aceita opções em colchetes. \\
\bottomrule
\end{longtable}

\section{Elementos pré e pós-textuais}

\begin{longtable}{>{\ttfamily\footnotesize\raggedright\arraybackslash}p{0.40\textwidth}p{0.22\textwidth}p{0.28\textwidth}}
\toprule
\normalfont\textbf{Comando} & \textbf{Resultado} & \textbf{Observações} \\
\midrule
\endhead
\string\undfcover & Capa & Usa as informações gerais. \\
\string\undffolhaderosto & Folha de rosto & Usa natureza e orientação. \\
\string\undfversofolhaderostoopcional & Ficha no verso & Só aparece com conteúdo configurado. \\
\string\undferrata & Errata & Só aparece quando preenchida. \\
\string\undffolhadeaprovacao & Folha de aprovação & Usa data e membros da banca. \\
\string\listadeilustracoes & Lista de figuras & Gerada a partir das legendas. \\
\string\listadefluxogramas & Lista própria & Exibida quando a contagem exige. \\
\string\listadetabelas & Lista de tabelas & Gerada a partir das legendas. \\
\string\undfsumario & Sumário & Exige novas passagens após mudanças. \\
\string\printundfreferences & Referências & Usa a base ligada por \string\addbibresource. \\
\string\undfapendice & Apêndice & Parâmetros: título e rótulo. \\
\string\undfanexo & Anexo & Parâmetros: título e rótulo. \\
\string\undfapendicepdf & Apêndice em PDF & Parâmetros: título e arquivo. \\
\string\undfanexarpdf & Anexo em PDF & Parâmetros: título e arquivo. \\
\bottomrule
\end{longtable}

\section{Citações e referências}

\begin{longtable}{>{\ttfamily\footnotesize\raggedright\arraybackslash}p{0.40\textwidth}p{0.22\textwidth}p{0.28\textwidth}}
\toprule
\normalfont\textbf{Comando} & \textbf{Resultado} & \textbf{Erro comum} \\
\midrule
\endhead
\string\citarparenteses & Citação parentética & Localizador opcional e chave obrigatória. \\
\string\citartexto & Autor integrado ao texto & Localizador opcional e chave obrigatória. \\
\string\citardireta & Citação curta parentética & Localizador, chave e texto. \\
\string\citardiretatexto & Citação curta narrativa & Localizador, chave e texto. \\
\string\citacaolonga & Bloco de citação longa & Localizador opcional e chave obrigatória. \\
\string\parenciteapud & Citação de citação & Autoria/ano originais e chave da consultada. \\
\string\textapud & Autoria integrada ao texto & Mesmos dados e localizadores opcionais. \\
\string\apudbib & Interface compatível anterior & Dois localizadores e duas chaves. \\
\string\fullref & Referência completa pontual & Parâmetro: chave. \\
\string\notarodape & Nota de rodapé & Parâmetro: texto. \\
\bottomrule
\end{longtable}

\section{Objetos e referências cruzadas}

\begin{longtable}{>{\ttfamily\footnotesize\raggedright\arraybackslash}p{0.43\textwidth}p{0.20\textwidth}p{0.27\textwidth}}
\toprule
\normalfont\textbf{Comando ou ambiente} & \textbf{Resultado} & \textbf{Observações} \\
\midrule
\endhead
\string\undfsecao & Seção rotulada & Parâmetros: título e rótulo. \\
\string\undfsubsecao & Subseção rotulada & Parâmetros: título e rótulo. \\
\string\equacao & Equação numerada & Rótulo opcional e expressão. \\
\string\undffigura & Figura & Largura, arquivo, legenda, rótulo e fonte. \\
\string\undffigurasladoalado... & Duas figuras independentes & Cada figura recebe seus próprios dados. \\
\string\undftabela & Tabela flutuante & Legenda, rótulo e fonte opcional. \\
\string\undffluxograma & Fluxograma & Legenda, rótulo e fonte opcional. \\
\string\undfpseudocodigo & Pseudocódigo comum & Assinatura, rótulo e título. \\
\string\undfpseudocodigolongo & Pseudocódigo multipágina & Mesmos parâmetros; permite quebra. \\
\string\figuraref & Chamada de figura & Parâmetro: rótulo sem \texttt{fig:}. \\
\string\tabelaref & Chamada de tabela & Parâmetro: rótulo sem \texttt{tab:}. \\
\string\equacaoref & Chamada de equação & Parâmetro: rótulo sem \texttt{eq:}. \\
\string\pseudocodigoref & Chamada de pseudocódigo & Parâmetro: rótulo sem \texttt{alg:}. \\
\bottomrule
\end{longtable}

\section{Validações}

\begin{longtable}{>{\ttfamily\footnotesize\raggedright\arraybackslash}p{0.40\textwidth}p{0.24\textwidth}p{0.26\textwidth}}
\toprule
\normalfont\textbf{Comando} & \textbf{Efeito} & \textbf{Observação} \\
\midrule
\endhead
\string\undfvalidateonepagecontents & Executa estimativas & Chame antes da capa. \\
\string\bloquearresumoslongos & Excesso vira erro & Opção de revisão, não padrão. \\
\string\avisarresumoslongos & Excesso permanece aviso & Comportamento padrão. \\
\string\undfdesabilitaravisosimagem & Suspende avisos de peso & Não altera a imagem. \\
\string\undfhabilitaravisosimagem & Reativa avisos & Ativo por padrão. \\
\bottomrule
\end{longtable}

\section{Exemplos e resultados esperados}

As duas tabelas seguintes se complementam. A primeira apresenta finalidade, parâmetros e uso básico; a segunda registra resultado esperado, erro comum e solução. Os comandos de apoio menos frequentes permanecem documentados nas tabelas temáticas anteriores.

\begin{landscape}
\begin{longtable}{>{\ttfamily\footnotesize\raggedright\arraybackslash}p{0.25\linewidth}p{0.19\linewidth}p{0.22\linewidth}>{\ttfamily\footnotesize\raggedright\arraybackslash}p{0.26\linewidth}}
\toprule
\normalfont\textbf{Comando} & \textbf{Finalidade} & \textbf{Parâmetros} & \normalfont\textbf{Uso básico} \\
\midrule
\endhead
\string\autor & Registrar autoria & Nome obrigatório & \string\autor\{Ana Lima\} \\
\string\citarparenteses & Inserir citação parentética & Localizador opcional e chave obrigatória & \string\citarparenteses\{silva2025\} \\
\string\undffigura & Inserir figura & Largura opcional, arquivo, título, rótulo e fonte opcional & \string\undffigura\{img\}\{Título\}\{id\} \\
\string\equacao & Inserir equação & Rótulo opcional e expressão obrigatória & \string\equacao[e]\{E=mc\string^2\} \\
\string\undfsecao & Criar seção rotulada & Título e rótulo obrigatórios & \string\undfsecao\{Objetivos\}\{obj\} \\
\string\bloquearresumoslongos & Ativar erro optativo de revisão & Nenhum & \string\bloquearresumoslongos \\
\bottomrule
\end{longtable}

\begin{longtable}{>{\ttfamily\footnotesize\raggedright\arraybackslash}p{0.25\linewidth}p{0.23\linewidth}p{0.21\linewidth}p{0.23\linewidth}}
\toprule
\normalfont\textbf{Comando} & \textbf{Resultado esperado} & \textbf{Erro comum} & \textbf{Solução} \\
\midrule
\endhead
\string\autor & Autoria nas páginas iniciais & Deixar o exemplo padrão & Substituir pelo nome real antes de compilar \\
\string\citarparenteses & Chamada autor-data vinculada à bibliografia & Usar chave inexistente & Conferir a chave no arquivo bibliográfico e executar o Biber \\
\string\undffigura & Figura numerada, com legenda, rótulo e fonte & Informar caminho incorreto & Conferir nome, extensão e caminho relativo do arquivo \\
\string\equacao & Equação numerada e referenciável & Repetir o rótulo & Usar um identificador único e recompilar \\
\string\undfsecao & Seção numerada e referenciável & Repetir o rótulo & Adotar rótulo único e descritivo \\
\string\bloquearresumoslongos & Excesso estimado tratado como erro & Ativar como comportamento comum de escrita & Reservar o comando para revisão controlada; o padrão oficial é aviso \\
\bottomrule
\end{longtable}
\end{landscape}


\chapter*{Referências Técnicas e Documentos Relacionados}
\addcontentsline{toc}{chapter}{Referências Técnicas e Documentos Relacionados}

\begin{itemize}
  \item \emph{Manual de Trabalhos de Conclusão de Curso da UnDF}, versão 1.0.0;
  \item documentação pública do projeto: \href{https://github.com/igormrqs/undftex-publico/blob/main/README.md}{README}, \href{https://github.com/igormrqs/undftex-publico/blob/main/docs/guia-de-uso.md}{guia de uso} e \href{https://github.com/igormrqs/undftex-publico/blob/main/docs/estrutura-do-projeto.md}{estrutura do projeto};
  \item documentação do LuaLaTeX, Biber, BibLaTeX e dos pacotes carregados pela classe;
  \item LaTeX Project Public License, versão 1.3c ou posterior;
  \item código-fonte da base estável do UnDF-TeX: \url{https://github.com/igormrqs/undftex-publico}.
\end{itemize}

\clearpage
\chapter*{Controle Documental}
\addcontentsline{toc}{chapter}{Controle Documental}

\noindent\begin{tabularx}{\textwidth}{@{}>{\bfseries\raggedright\arraybackslash}p{0.31\textwidth}Y@{}}
\toprule
Título & Guia de Uso do UnDF-TeX \\
Versão do manual & 0.3.1-preliminar \\
Compatibilidade & UnDF-TeX 0.1.0 \\
Situação & Minuta para validação institucional \\
Data & 7 de agosto de 2026 \\
Escopo & Preparação, edição, compilação e verificação técnica do template \\
\bottomrule
\end{tabularx}


\end{document}
